\section{Conversion also changes the accounting question}
\label{sim:conversion-control-scope}

The economic financing model specifies who receives cash and who chooses the project. Accounting for those events requires additional distinctions. In the supplied contingent-conversion extension, a creditor receives 120 shares alongside 40 existing shares. The resulting 75 per cent voting fraction explains the extension's assumed economic control rule. It does not by itself establish the date control begins for accounting, the holder's reporting perimeter, or the amount to recognise in its accounts.

The following analysis uses Australian standards as a conditional design case. It does not select a reporting basis for an actual company. The retained editions and their application periods matter: the inspected AASB 10 compilation covers annual periods beginning in 2026 before 2027, and the inspected AASB 3 compilation covers starts from 1 July 2026 before 2027. The synthetic issue and conversion occur in 2027 and 2028, so a complete applicable edition and entity-scope chain remains necessary. The source paragraphs identify the questions that must be answered, while the missing facts remain unresolved.

\subsection{Control can begin before the shares arrive}

Under AASB 10, control requires power over the investee, exposure or rights to variable returns, and the ability to use that power to affect the returns. Power concerns existing rights giving the current ability to direct the activities that significantly affect those returns. Majority votes support that conclusion in the circumstances specified by the standard, but the votes must be substantive and contrary rights or arrangements can matter.\citep[paras.~6--18, B22--B28, B35--B37]{aasb10control2026}

This is a question about rights and relevant decisions, not just a percentage. A lender's protective covenant can restrict a borrower without giving the lender power over its activities. Conversely, an apparently contingent decision right is not automatically protective. If the activities that significantly affect returns arise only after a particular event, the right to direct those activities can matter before the event occurs.\citep[paras.~B26--B28, B53]{aasb10control2026}

Potential voting rights introduce the same timing issue into a convertible instrument. AASB 10 requires consideration of substantive potential rights, including those arising from convertibles, options and forwards. Whether rights are substantive depends partly on whether they can be exercised when decisions about the relevant activities need to be made.\citep[paras.~B24, B47--B50]{aasb10control2026} It follows that the simulator needs a control assessment before conversion as well as after it. Setting the first control date equal to the share-issue timestamp would conceal this question.

The acquisition date for an in-scope business combination is the date control is obtained. It can precede or follow the legal closing date when the relevant facts support that timing. Consolidation similarly begins when control is obtained.\citep[para.~20]{aasb10control2026}\citep[paras.~8--9]{aasb3scope2026} Thus the economic conversion, the rights assessment and the accounting acquisition date require distinct records linked by evidence. A date cannot be supplied merely to make the annual schedule continuous.

\subsection{The holder's investment and the group's assets are different reports}

In separate financial statements, AASB 127 permits investments in subsidiaries, joint ventures and associates to be accounted for at cost, under AASB 9, or using the equity method, subject to its requirements and consistent treatment within each investment category. Certain required fair-value treatments carry through to the separate statements.\citep[paras.~3--10, 11A]{aasb127scope2022} Choosing a cost policy does not provide the initial cost amount. The earlier price paid for a convertible claim cannot silently become the post-conversion measurement without the relevant transaction and measurement analysis.

For consolidated reporting, a parent normally consolidates controlled subsidiaries, with specified exemptions and exceptions. An investment entity generally measures qualifying subsidiary investments at fair value through profit or loss, while a non-investment-entity subsidiary providing services related to its parent's investment activities can fall within the consolidation exception to that exception.\citep[paras.~4--4B, 27--33]{aasb10control2026} The investor's entity status, the subsidiary's activities and the complete exemption conditions therefore belong among the required facts. None is determined by the 75 per cent fraction.

Consolidation combines the relevant assets, liabilities, income, expenses and cash flows, eliminates the parent's investment against the relevant subsidiary equity, and removes intragroup transactions and balances. It also requires consistent policies and appropriate reporting dates.\citep[paras.~B86--B93]{aasb10control2026} Potential rights used in a control assessment do not automatically set the proportion allocated to non-controlling interests: the standard separately addresses existing ownership and instruments that already provide access to ownership returns.\citep[paras.~B89--B91]{aasb10control2026}

This creates a concrete implementation consequence. The existing selected group-period reference assumes the same supplied control edges and membership at opening and closing. A company acquired during that period violates the assumption. A correct continuation must distinguish the period before control, the acquisition measurements and the period consolidated afterward, rather than force the new subsidiary into the opening group or omit its acquisition effects.

\subsection{A controlled project is not automatically a business combination}

Obtaining control does not settle whether the acquired set is a business. AASB 3 distinguishes an in-scope business combination from acquiring an asset or group of assets, a common-control combination, and specified investment-entity acquisitions. An asset-group acquisition under its scope paragraph does not create goodwill.\citep[paras.~2--3]{aasb3scope2026}

The optional concentration test can establish that an acquired set is not a business when substantially all of the gross-asset fair value is concentrated as specified. Failing that test does not establish that the set is a business; the further assessment remains necessary. The underlying analysis requires inputs and a substantive process, with different criteria depending on whether outputs are present.\citep[paras.~B7A--B12D]{aasb3scope2026} A project described by a terminal cash payoff supplies neither the workforce and process evidence nor the fair-value composition needed to decide the question.

For an in-scope business combination, the acquisition method requires identifying the acquirer and control date, recognising and measuring the identifiable assets, liabilities and non-controlling interest, and determining goodwill or a bargain-purchase gain. The goodwill calculation uses the specified acquisition-date components, including any relevant previously held interest.\citep[paras.~4--10, 32--33]{aasb3scope2026} The synthetic funding price, a 5 per cent pricing discount and a balanced four-party journal do not supply these missing measurements. They remain separate economic evidence.

\subsection{Keep dated authority and missing facts visible}

The inspected AASB 18 consequential amendments change specified AASB 3 presentation and disclosure references, add title-related footnotes in AASB 10 and AASB 127, and alter accompanying AASB 10 guidance. They do not in that instrument rewrite the core control, business-definition or separate-policy paragraphs used here.\citep[Appendix D, pp.~57, 70, 101--102]{aasb18scope2024} That bounded observation cannot establish that no other amendment applies.

AASB 2024-4 defers specified sale-or-contribution amendments involving an investor and its associate or joint venture until annual periods beginning in 2028, with earlier application permitted. It is not a general acquisition-of-control rule.\citep[paras.~6--11]{aasb2024controltiming} AASB 2026-3 changes a downstream associate/joint-venture fair-value election, while AASB 2026-4 contains entity-specific presentation and application changes and references AASB 2026-2.\citep[p.~5]{aasb2026associate}\citep[paras.~1--3]{aasb2026presentation} Completing that entity and edition chain is still required before actual 2027 or 2028 admission.

The conditional dependency reference consequently records supplied, dated assessments of rights, control, reporting scope and acquisition type. Unknown assessments stay unknown; missing measurements stay named requirements. It distinguishes a separate cost-policy choice from evidence of cost, and a positive supplied control assessment from actual source-backed admission. Its current results authorize no posting. This is the useful conclusion for scenario design: unresolved accounting alternatives must remain linked to the rights, dates, entities and measurements that could distinguish them, instead of being collapsed into whichever journal first balances.\footnote{Retained official sources, paragraph and physical-page anchors, the conditional analysis and refusal cases are in \path{data/accounting-master/conversion-control-scope-20261026/}. The selected dependency reference checks 21 supplied scenarios and 12 invalid-input cases; those checks establish reference behavior, not actual accounting correctness or complete standards coverage.}
